% ============================================================================
% MANUAL DE USUARIO — SADUTO
% Sistema de Analítica de Datos Utilizando Técnicas Open Source Intelligence
% Caso: Vicerrectorado de Grado — Escuela Militar de Ingeniería
% ============================================================================
\documentclass[12pt,a4paper]{article}

% ── Paquetes ──────────────────────────────────────────────────────────────
\usepackage[utf8]{inputenc}
\usepackage[spanish]{babel}
\usepackage[T1]{fontenc}
\usepackage{lmodern}
\usepackage[margin=2.5cm]{geometry}
\usepackage{graphicx}
\usepackage{xcolor}
\usepackage{titlesec}
\usepackage{fancyhdr}
\usepackage{hyperref}
\usepackage{enumitem}
\usepackage{booktabs}
\usepackage{longtable}
\usepackage{tabularx}
\usepackage{float}
\usepackage{tcolorbox}
\usepackage{fontawesome5}
\usepackage{array}
\usepackage{pifont}

% ── Colores institucionales ───────────────────────────────────────────────
\definecolor{emiblue}{HTML}{0D47A1}
\definecolor{emiyellow}{HTML}{FFD700}
\definecolor{emidark}{HTML}{002171}
\definecolor{emigray}{HTML}{F5F5F5}
\definecolor{emilightblue}{HTML}{E3F2FD}

% ── Configuración de hipervínculos ────────────────────────────────────────
\hypersetup{
  colorlinks=true,
  linkcolor=emiblue,
  urlcolor=emiblue,
  citecolor=emiblue,
  pdftitle={Manual de Usuario — SADUTO},
  pdfauthor={Alvaro Santiago Encinas Flores},
}

% ── Formato de títulos ────────────────────────────────────────────────────
\titleformat{\section}{\LARGE\bfseries\color{emiblue}}{}{0em}{}[\titlerule]
\titleformat{\subsection}{\Large\bfseries\color{emidark}}{}{0em}{}
\titleformat{\subsubsection}{\large\bfseries\color{emiblue}}{}{0em}{}

% ── Encabezado y pie de página ────────────────────────────────────────────
\pagestyle{fancy}
\fancyhf{}
\fancyhead[L]{\small\color{emiblue}\textbf{SADUTO} — Manual de Usuario}
\fancyhead[R]{\small\color{gray}v1.0}
\fancyfoot[C]{\thepage}
\fancyfoot[R]{\small\color{gray}EMI Bolivia — 2026}
\renewcommand{\headrulewidth}{1pt}
\renewcommand{\headrule}{\hbox to\headwidth{\color{emiblue}\leaders\hrule height \headrulewidth\hfill}}

% ── Cajas de información ─────────────────────────────────────────────────
\newtcolorbox{infobox}[1][]{
  colback=emilightblue,
  colframe=emiblue,
  fonttitle=\bfseries,
  title=#1,
  rounded corners,
  boxrule=0.5pt
}

\newtcolorbox{warningbox}[1][]{
  colback=yellow!10,
  colframe=emiyellow!80!black,
  fonttitle=\bfseries,
  title=#1,
  rounded corners,
  boxrule=0.5pt
}

\newtcolorbox{tipbox}[1][]{
  colback=green!5,
  colframe=green!60!black,
  fonttitle=\bfseries,
  title=#1,
  rounded corners,
  boxrule=0.5pt
}

% ══════════════════════════════════════════════════════════════════════════
\begin{document}

% ── Portada ───────────────────────────────────────────────────────────────
\begin{titlepage}
\centering
\vspace*{2cm}

{\Huge\bfseries\color{emiblue}SADUTO\par}
\vspace{0.3cm}
{\large\color{emidark}Sistema de Analítica de Datos Utilizando\\Técnicas Open Source Intelligence\\Caso: Vicerrectorado de Grado\par}

\vspace{1cm}
\noindent\makebox[\textwidth]{\color{emiblue}\rule{0.6\textwidth}{2pt}}
\vspace{1cm}

{\LARGE\bfseries Manual de Usuario\par}
\vspace{0.3cm}
{\large Guía completa para usuarios finales y administradores\par}

\vfill

\begin{tabular}{ll}
\textbf{Autor:}         & Alvaro Santiago Encinas Flores \\
\textbf{Carrera:}       & Ing.\ de Sistemas — EMI 2026 \\
\textbf{Institución:}   & Escuela Militar de Ingeniería (EMI) \\
\textbf{Dependencia:}   & Vicerrectorado de Grado \\
\textbf{Unidad:}        & Unidad de Evaluación y Bienestar Universitario (UEBU) \\
\textbf{Versión:}       & 1.0 \\
\textbf{Fecha:}         & Febrero 2026 \\
\textbf{Clasificación:} & Uso Interno \\
\end{tabular}

\vspace{1cm}
\noindent\makebox[\textwidth]{\color{emiyellow}\rule{0.4\textwidth}{3pt}}
\end{titlepage}

% ── Tabla de contenidos ───────────────────────────────────────────────────
\tableofcontents
\newpage

% ══════════════════════════════════════════════════════════════════════════
% PARTE I — MANUAL DE USUARIO FINAL
% ══════════════════════════════════════════════════════════════════════════

\section{Introducción}

\subsection{Propósito del documento}
Este manual proporciona instrucciones detalladas para el uso del sistema \textbf{SADUTO} (Sistema de Analítica de Datos Utilizando Técnicas Open Source Intelligence — Caso: Vicerrectorado de Grado). Está dirigido a:

\begin{itemize}[leftmargin=2cm]
  \item[\faUserShield] \textbf{Administradores del sistema} — Gestión técnica y configuración.
  \item[\faUserTie] \textbf{Vicerrector de Grado} — Consulta de reportes ejecutivos y KPIs estratégicos.
  \item[\faUsers] \textbf{Coordinadores y Analistas UEBU} — Operación diaria, análisis de datos e interpretación de dashboards.
\end{itemize}

\subsection{Descripción del sistema}
SADUTO es una plataforma web que recolecta, procesa y analiza datos de fuentes abiertas (OSINT) sobre la percepción institucional de la EMI en redes sociales y medios digitales. El sistema utiliza inteligencia artificial para:

\begin{enumerate}
  \item Extraer publicaciones y comentarios de Facebook y TikTok.
  \item Clasificar el sentimiento (positivo, neutral, negativo) mediante modelos de Deep Learning.
  \item Detectar temas, tendencias y anomalías con técnicas NLP/ML.
  \item Generar reportes ejecutivos en PDF y Excel.
  \item Emitir alertas automáticas ante situaciones críticas.
\end{enumerate}

\subsection{Requisitos del navegador}
\begin{table}[H]
\centering
\begin{tabular}{lcc}
\toprule
\textbf{Navegador} & \textbf{Versión mínima} & \textbf{Recomendado} \\
\midrule
Google Chrome & 90+ & \ding{51} \\
Mozilla Firefox & 88+ & \ding{51} \\
Microsoft Edge & 90+ & \ding{51} \\
Safari & 14+ & — \\
\bottomrule
\end{tabular}
\caption{Navegadores compatibles}
\end{table}

% ──────────────────────────────────────────────────────────────────────────
\section{Acceso al sistema}

\subsection{Inicio de sesión}
\begin{enumerate}
  \item Abra su navegador y acceda a la URL del sistema: \texttt{http://[servidor]:3000/login}
  \item Ingrese su \textbf{correo electrónico institucional} (ej.\ \texttt{admin@emi.edu.bo}).
  \item Ingrese su \textbf{contraseña} asignada por el administrador.
  \item Haga clic en el botón \textbf{``Iniciar Sesión''}.
\end{enumerate}

\begin{infobox}[\faInfoCircle\ Nota]
La pantalla de login muestra el logo institucional SADUTO sobre un fondo con overlay azul. Los campos de texto están resaltados con el color institucional azul (\texttt{\#0D47A1}).
\end{infobox}

\subsection{Roles de usuario}
El sistema maneja tres perfiles de acceso con diferentes niveles de permisos:

\begin{table}[H]
\centering
\begin{tabularx}{\textwidth}{l l X}
\toprule
\textbf{Rol} & \textbf{Código} & \textbf{Permisos} \\
\midrule
Administrador & \texttt{administrador} & Acceso completo: CRUD de usuarios, configuración de alertas, todos los dashboards, logs del sistema. \\
\addlinespace
Vicerrector & \texttt{vicerrector} & Consulta de dashboards, reportes ejecutivos, KPIs estratégicos. Sin acceso a gestión de usuarios. \\
\addlinespace
Analista UEBU & \texttt{uebu} & Operación de dashboards, extracción de datos, generación de reportes, configuración de alertas propias. \\
\bottomrule
\end{tabularx}
\caption{Roles y permisos del sistema}
\end{table}

\subsection{Recuperación de contraseña}
Si olvidó su contraseña, contacte al administrador del sistema para que le asigne una nueva desde el módulo de Gestión de Usuarios.

% ──────────────────────────────────────────────────────────────────────────
\section{Navegación general}

\subsection{Estructura de la interfaz}
La interfaz de SADUTO se compone de tres áreas principales:

\begin{enumerate}
  \item \textbf{Barra superior (Header)}: Muestra el logo SADUTO, el título de la sección actual, controles de tema claro/oscuro, notificaciones y menú de perfil de usuario.
  \item \textbf{Panel lateral (Sidebar)}: Menú de navegación con las siguientes secciones:
  \begin{itemize}
    \item Inteligencia OSINT
    \item Posts y Comentarios
    \item Análisis AI (submenú: Sentimientos, Reputación, Alertas, Benchmarking)
    \item IA / ML / NLP
    \item Evaluación del Sistema
    \item Gestión de Usuarios
    \item Configuración
  \end{itemize}
  \item \textbf{Área principal}: Espacio de contenido donde se muestran los dashboards, tablas y formularios.
\end{enumerate}

\subsection{Tema visual}
El sistema soporta dos modos de visualización:
\begin{itemize}
  \item \textbf{Modo claro}: Fondo claro con texto oscuro (recomendado para uso diurno).
  \item \textbf{Modo oscuro}: Fondo oscuro con texto claro (reduce fatiga visual en uso nocturno).
\end{itemize}
Puede alternar entre ambos modos usando el ícono de sol/luna en la barra superior.

% ──────────────────────────────────────────────────────────────────────────
\section{Módulo: Inteligencia OSINT}

\subsection{Descripción}
Este dashboard presenta una visión consolidada de toda la inteligencia recolectada de fuentes abiertas. Incluye:

\begin{itemize}
  \item \textbf{KPIs principales}: Total de fuentes activas, datos recolectados, patrones identificados y noticias monitoreadas.
  \item \textbf{Noticias relevantes}: Feed de noticias sobre la EMI extraídas de medios digitales, con indicador de relevancia y sentimiento.
  \item \textbf{Tendencias Google}: Gráficos de interés de búsqueda de términos relacionados con la EMI en Bolivia.
  \item \textbf{Patrones identificados}: Patrones de comportamiento detectados por los algoritmos de IA.
\end{itemize}

\subsection{Interpretación de KPIs OSINT}
\begin{table}[H]
\centering
\begin{tabularx}{\textwidth}{l X l}
\toprule
\textbf{KPI} & \textbf{Descripción} & \textbf{Valor ideal} \\
\midrule
Fuentes activas & Número de plataformas monitoreadas & $\geq 3$ \\
Datos recolectados & Volumen total de publicaciones/comentarios & Creciente \\
Patrones detectados & Conductas repetitivas identificadas & Variable \\
Noticias & Artículos sobre la EMI en medios & Monitoreo \\
\bottomrule
\end{tabularx}
\caption{KPIs del módulo OSINT}
\end{table}

% ──────────────────────────────────────────────────────────────────────────
\section{Módulo: Posts y Comentarios}

\subsection{Descripción}
Permite gestionar y visualizar las publicaciones (posts) extraídas de las fuentes OSINT configuradas (Facebook y TikTok), así como sus comentarios asociados.

\subsection{Funcionalidades}
\begin{enumerate}
  \item \textbf{Lista de fuentes}: Panel izquierdo con las fuentes OSINT configuradas.
  \item \textbf{Visualización de posts}: Lista central con títulos, fechas, engagement y sentimiento.
  \item \textbf{Detalle de comentarios}: Al seleccionar un post, se despliegan sus comentarios con análisis de sentimiento individual.
  \item \textbf{Extracción de comentarios TikTok}: Botón para iniciar scraping en tiempo real de comentarios de TikTok.
\end{enumerate}

\begin{warningbox}[\faExclamationTriangle\ Precaución]
La extracción de comentarios de TikTok puede tardar varios minutos dependiendo de la cantidad de comentarios. No cierre el navegador durante este proceso.
\end{warningbox}

% ──────────────────────────────────────────────────────────────────────────
\section{Módulo: Análisis de Sentimientos}

\subsection{Descripción}
Dashboard central de inteligencia artificial que muestra los resultados del análisis de sentimiento sobre todos los datos recolectados.

\subsection{Componentes del dashboard}
\begin{enumerate}
  \item \textbf{KPIs de sentimiento}:
  \begin{itemize}
    \item \textbf{Índice de Satisfacción General (ISG)}: Porcentaje calculado como:
    $$ISG = \frac{\text{Publicaciones positivas}}{\text{Total publicaciones}} \times 100$$
    \item \textbf{Ratio positivo/negativo}: Proporción entre comentarios positivos y negativos.
    \item \textbf{Confianza promedio del modelo}: Nivel de certeza del clasificador de sentimiento.
  \end{itemize}
  \item \textbf{Gráfico de distribución}: Gráfico circular (pie chart) con la proporción de sentimientos positivo, neutral y negativo.
  \item \textbf{Gráfico de tendencia temporal}: Línea de tiempo mostrando la evolución del sentimiento.
  \item \textbf{Nube de palabras}: Palabras más frecuentes en los comentarios, ponderadas por TF-IDF.
  \item \textbf{Posts recientes}: Lista de las últimas publicaciones con su clasificación de sentimiento.
\end{enumerate}

\subsection{Interpretación de colores de sentimiento}
\begin{table}[H]
\centering
\begin{tabular}{lcl}
\toprule
\textbf{Sentimiento} & \textbf{Color} & \textbf{Significado} \\
\midrule
Positivo & \textcolor{green!60!black}{\rule{1cm}{0.4cm}} & La publicación expresa satisfacción o aprobación. \\
Neutral  & \textcolor{gray}{\rule{1cm}{0.4cm}} & La publicación es informativa, sin carga emocional clara. \\
Negativo & \textcolor{red!70!black}{\rule{1cm}{0.4cm}} & La publicación expresa insatisfacción o queja. \\
\bottomrule
\end{tabular}
\caption{Código de colores para sentimiento}
\end{table}

% ──────────────────────────────────────────────────────────────────────────
\section{Módulo: Reputación}

\subsection{Descripción}
Proporciona una visión integral de la reputación digital de la EMI, comparando el desempeño institucional frente a competidores y mostrando los temas más relevantes.

\subsection{Componentes}
\begin{itemize}
  \item \textbf{Score de Reputación}: Indicador compuesto (0--100) que pondera sentimiento, volumen y engagement.
  \item \textbf{Temas de reputación}: Categorías temáticas identificadas (infraestructura, académico, becas, etc.).
  \item \textbf{Competidores}: Comparativa de menciones y sentimiento vs.\ otras universidades bolivianas.
\end{itemize}

% ──────────────────────────────────────────────────────────────────────────
\section{Módulo: Alertas y Anomalías}

\subsection{Descripción}
Gestiona las alertas generadas automáticamente por el sistema cuando se detectan situaciones que requieren atención.

\subsection{Tipos de alertas}
\begin{table}[H]
\centering
\begin{tabularx}{\textwidth}{l l X}
\toprule
\textbf{Severidad} & \textbf{Color} & \textbf{Criterio de activación} \\
\midrule
Crítica & \textcolor{red}{\textbf{Rojo}} & Sentimiento negativo $>70\%$ con alta confianza, picos de quejas. \\
Alta & \textcolor{orange}{\textbf{Naranja}} & Incremento significativo de menciones negativas. \\
Media & \textcolor{yellow!80!black}{\textbf{Amarillo}} & Tendencias descendentes moderadas en sentimiento. \\
Baja & \textcolor{blue}{\textbf{Azul}} & Cambios menores que requieren monitoreo. \\
\bottomrule
\end{tabularx}
\caption{Niveles de severidad de alertas}
\end{table}

\subsection{Gestión de alertas}
\begin{enumerate}
  \item Las alertas aparecen en el dashboard con su severidad, título y descripción.
  \item Haga clic en una alerta para ver su detalle completo.
  \item Use el botón \textbf{``Resolver''} para marcar una alerta como resuelta, incluyendo notas de resolución.
  \item Las alertas resueltas se mueven al historial.
\end{enumerate}

% ──────────────────────────────────────────────────────────────────────────
\section{Módulo: Benchmarking}

\subsection{Descripción}
Compara métricas de percepción digital de las diferentes carreras de la EMI, permitiendo identificar las que tienen mejor y peor desempeño en redes sociales.

\subsection{Métricas comparativas}
\begin{itemize}
  \item \textbf{Menciones}: Volumen de publicaciones que mencionan cada carrera.
  \item \textbf{Sentimiento}: Porcentaje de percepción positiva por carrera.
  \item \textbf{Engagement}: Nivel de interacción (likes, comentarios, compartidos) normalizado.
\end{itemize}

% ──────────────────────────────────────────────────────────────────────────
\section{Módulo: IA / ML / NLP}

\subsection{Descripción}
Dashboard técnico que muestra los resultados de los algoritmos de Procesamiento de Lenguaje Natural y Machine Learning.

\subsection{Componentes}
\begin{enumerate}
  \item \textbf{Palabras clave (TF-IDF)}: Las 50 palabras más relevantes extraídas por el algoritmo Term Frequency -- Inverse Document Frequency.
  \item \textbf{Modelado de tópicos (LDA)}: Agrupaciones temáticas generadas por Latent Dirichlet Allocation.
  \item \textbf{Clustering (K-Means)}: Agrupación de documentos por similitud semántica, visualizado con etiquetas y sentimiento predominante.
  \item \textbf{Entidades nombradas (NER)}: Personas, organizaciones y lugares mencionados en los datos.
  \item \textbf{Clasificación temática}: Distribución de contenido por categorías (académico, infraestructura, becas, etc.).
  \item \textbf{Resumen ejecutivo}: Texto generado automáticamente que sintetiza los hallazgos principales.
\end{enumerate}

\begin{tipbox}[\faLightbulb\ Consejo]
Los resultados de NLP se actualizan cada vez que se ejecuta el pipeline de análisis. Consulte con el administrador la frecuencia de actualización configurada.
\end{tipbox}

% ──────────────────────────────────────────────────────────────────────────
\section{Módulo: Evaluación del Sistema}

\subsection{Descripción}
Presenta métricas de rendimiento y calidad del sistema SADUTO, alineadas con los objetivos específicos de la investigación.

\subsection{Categorías de evaluación}
\begin{table}[H]
\centering
\begin{tabularx}{\textwidth}{l c X}
\toprule
\textbf{Categoría} & \textbf{Peso} & \textbf{Qué mide} \\
\midrule
Recolección & 20\% & Volumen, diversidad de fuentes, calidad de datos, cobertura temporal. \\
Sentimiento & 25\% & Cobertura del análisis, balance, confianza del modelo, uso de Deep Learning. \\
NLP/ML & 25\% & Keywords, Topic Modeling, Clustering, NER, clasificación temática. \\
Rendimiento & 15\% & Tiempos de respuesta de la API ($<2$ seg). \\
Completitud & 15\% & Estructura completa de la base de datos. \\
\bottomrule
\end{tabularx}
\caption{Categorías de evaluación del sistema}
\end{table}

\subsection{Objetivos específicos evaluados}
\begin{description}
  \item[OE1] Implementar recolección multi-fuente OSINT (Facebook, TikTok, Google News).
  \item[OE2] Desarrollar dashboards interactivos (8 dashboards implementados).
  \item[OE3] Aplicar AI/ML/NLP para análisis de sentimiento y detección de patrones.
  \item[OE4] Evaluar la calidad del sistema con métricas cuantitativas ($\geq 90/100$).
\end{description}

% ──────────────────────────────────────────────────────────────────────────
\section{Módulo: Gestión de Usuarios}

\subsection{Descripción}
\textit{(Solo disponible para rol Administrador)}. Permite realizar operaciones CRUD sobre las cuentas de usuario del sistema.

\subsection{Operaciones disponibles}
\begin{enumerate}
  \item \textbf{Listar usuarios}: Tabla con todos los usuarios, su rol, estado y último acceso.
  \item \textbf{Crear usuario}: Formulario para agregar nuevos usuarios con asignación de rol.
  \item \textbf{Editar usuario}: Modificar nombre, correo, rol y estado (activo/inactivo).
  \item \textbf{Eliminar usuario}: Desactivar una cuenta de usuario.
  \item \textbf{Buscar}: Filtro de búsqueda por nombre, email o rol.
\end{enumerate}

\begin{infobox}[\faInfoCircle\ Seguridad]
Las contraseñas se almacenan con hash SHA-256. La contraseña por defecto para nuevos usuarios es \texttt{emi2024} y debe cambiarse en el primer inicio de sesión.
\end{infobox}

% ──────────────────────────────────────────────────────────────────────────
\section{Módulo: Configuración}

\subsection{Configuración de alertas}
Permite definir reglas personalizadas para la generación de alertas automáticas:
\begin{itemize}
  \item \textbf{Nombre de la regla}: Identificador descriptivo.
  \item \textbf{Tipo de alerta}: Sentimiento negativo, anomalía de volumen, etc.
  \item \textbf{Umbral de valor}: Porcentaje que activa la alerta (ej.\ 0.7 = 70\%).
  \item \textbf{Severidad mínima}: Nivel mínimo para generar la alerta.
  \item \textbf{Notificación por email}: Opción de envío automático.
\end{itemize}

\subsection{Perfil de usuario}
Cada usuario puede actualizar desde esta sección:
\begin{itemize}
  \item Su nombre completo y cargo.
  \item Su contraseña (requiere la contraseña actual).
\end{itemize}

% ──────────────────────────────────────────────────────────────────────────
\section{Centro de Reportes}

\subsection{Generación de reportes}
El sistema permite generar dos tipos de reportes:

\subsubsection{Reportes PDF}
\begin{table}[H]
\centering
\begin{tabularx}{\textwidth}{l X l}
\toprule
\textbf{Tipo} & \textbf{Contenido} & \textbf{Páginas} \\
\midrule
Reporte Ejecutivo & Resumen semanal con KPIs, gráficos de tendencias y recomendaciones. & 8--12 \\
Reporte de Alertas & Detalle de alertas críticas con timeline y acciones recomendadas. & 4--6 \\
Anuario Estadístico & Análisis semestral completo con todos los indicadores. & 30--50 \\
Reporte por Carrera & Análisis específico de una carrera con comparativas. & 10--15 \\
\bottomrule
\end{tabularx}
\caption{Tipos de reportes PDF}
\end{table}

\subsubsection{Reportes Excel}
\begin{table}[H]
\centering
\begin{tabularx}{\textwidth}{l X}
\toprule
\textbf{Tipo} & \textbf{Contenido} \\
\midrule
Dataset de Sentimientos & Datos completos con múltiples hojas de análisis y gráficos. \\
Tabla Pivote & Análisis agregado por carrera, fuente o periodo. \\
Reporte de Anomalías & Detección de patrones anómalos con alertas. \\
Reporte Combinado & Todas las métricas consolidadas en un solo archivo. \\
\bottomrule
\end{tabularx}
\caption{Tipos de reportes Excel}
\end{table}

\subsection{Historial de reportes}
Todos los reportes generados se almacenan y pueden:
\begin{itemize}
  \item Descargarse nuevamente desde el historial.
  \item Filtrarse por tipo, formato y fecha.
  \item Eliminarse si ya no son necesarios.
\end{itemize}

\subsection{Reportes programados}
Los administradores pueden configurar la generación automática de reportes:
\begin{itemize}
  \item \textbf{Frecuencia}: Diaria, semanal o mensual.
  \item \textbf{Distribución}: Envío automático por email a destinatarios configurados.
  \item \textbf{Formato}: PDF o Excel según configuración.
\end{itemize}

% ══════════════════════════════════════════════════════════════════════════
% PARTE II — MANUAL DE OPERACIONES / ADMINISTRADOR
% ══════════════════════════════════════════════════════════════════════════
\newpage
\section{Manual de Operaciones (Administrador)}

\subsection{Arquitectura del sistema}
\begin{table}[H]
\centering
\begin{tabularx}{\textwidth}{l l X}
\toprule
\textbf{Componente} & \textbf{Tecnología} & \textbf{Descripción} \\
\midrule
Backend API & Flask (Python 3.11) & API REST en puerto 5001, $\sim$3.350 líneas. \\
Frontend & React 18 + TypeScript + MUI 5 & Aplicación SPA en puerto 3000 con Vite. \\
Base de datos & SQLite 3 & Archivo \texttt{data/osint\_emi.db}, 23 tablas. \\
NLP/ML & NLTK, scikit-learn, transformers & Pipeline de análisis con modelos pre-entrenados. \\
Scraping & Selenium, requests & Extracción de Facebook y TikTok. \\
Reportes & ReportLab, openpyxl & Generación de PDF y Excel. \\
\bottomrule
\end{tabularx}
\caption{Stack tecnológico del sistema}
\end{table}

\subsection{Instalación desde cero}

\subsubsection{Prerrequisitos}
\begin{itemize}
  \item Python 3.11 o superior
  \item Node.js 18+ y npm
  \item Git
  \item Google Chrome (para Selenium)
\end{itemize}

\subsubsection{Paso 1: Clonar el repositorio}
\begin{verbatim}
git clone [URL_REPOSITORIO]
cd osint_vicerrectorado
\end{verbatim}

\subsubsection{Paso 2: Configurar el backend}
\begin{verbatim}
python3 -m venv venv
source venv/bin/activate        # macOS/Linux
venv\Scripts\activate           # Windows
pip install -r requirements.txt
\end{verbatim}

\subsubsection{Paso 3: Inicializar la base de datos}
\begin{verbatim}
python crear_bd.py
python seed_data.py     # Datos iniciales (fuentes, usuarios)
\end{verbatim}

\subsubsection{Paso 4: Configurar el frontend}
\begin{verbatim}
cd frontend
npm install
\end{verbatim}

\subsubsection{Paso 5: Iniciar los servicios}
\begin{verbatim}
# Terminal 1 — Backend
cd osint_vicerrectorado
./venv/bin/python api_real.py

# Terminal 2 — Frontend
cd osint_vicerrectorado/frontend
npm run dev
\end{verbatim}

\begin{infobox}[\faInfoCircle\ Verificación]
Acceda a \texttt{http://localhost:3000/login} para verificar que la instalación fue exitosa. Use las credenciales: \texttt{admin@emi.edu.bo} / \texttt{emi2024}.
\end{infobox}

\subsection{Variables de entorno}
\begin{table}[H]
\centering
\begin{tabularx}{\textwidth}{l l X}
\toprule
\textbf{Variable} & \textbf{Valor por defecto} & \textbf{Descripción} \\
\midrule
\texttt{FLASK\_PORT} & 5001 & Puerto de la API Flask. \\
\texttt{VITE\_API\_URL} & http://localhost:5001 & URL del backend para el frontend. \\
\texttt{DB\_PATH} & data/osint\_emi.db & Ruta de la base de datos SQLite. \\
\texttt{SECRET\_KEY} & osint-emi-secret-2024 & Clave secreta para tokens JWT. \\
\texttt{DEBUG} & True & Modo debug (False en producción). \\
\bottomrule
\end{tabularx}
\caption{Variables de entorno del sistema}
\end{table}

\subsection{Troubleshooting}

\begin{longtable}{p{3.5cm} p{3cm} p{7cm}}
\toprule
\textbf{Error} & \textbf{Código} & \textbf{Solución} \\
\midrule
\endhead

Vite chunk not found & ERR\_FILE\_NOT\_FOUND & Ejecutar: \texttt{rm -rf node\_modules/.vite \&\& npm run dev} \\
\addlinespace

\texttt{dispatcher.useContext} is null & TypeError & Limpiar caché de Vite y hacer hard refresh (Cmd+Shift+R). \\
\addlinespace

API no responde & Connection refused & Verificar que el backend está corriendo: \texttt{curl http://localhost:5001/api/health} \\
\addlinespace

Error de autenticación & 401 Unauthorized & Verificar token en header: \texttt{Authorization: Bearer <token>} \\
\addlinespace

Base de datos bloqueada & SQLITE\_BUSY & Cerrar otros procesos que accedan a la BD. Reiniciar la API. \\
\addlinespace

Selenium no encuentra Chrome & WebDriverError & Instalar Chrome y chromedriver compatible. \\
\addlinespace

Módulos Python no encontrados & ModuleNotFoundError & Activar el entorno virtual: \texttt{source venv/bin/activate} \\

\bottomrule
\caption{Guía de troubleshooting}
\end{longtable}

\subsection{Backup y restauración}

\subsubsection{Backup de la base de datos}
\begin{verbatim}
# Crear backup con timestamp
cp data/osint_emi.db data/backup/osint_emi_$(date +%Y%m%d_%H%M%S).db

# Backup programado (agregar a crontab)
0 2 * * * cp /ruta/data/osint_emi.db /ruta/data/backup/osint_emi_$(date +\%Y\%m\%d).db
\end{verbatim}

\subsubsection{Restauración}
\begin{verbatim}
# Detener el servicio API
pkill -f "python.*api_real"

# Restaurar desde backup
cp data/backup/osint_emi_20250223.db data/osint_emi.db

# Reiniciar la API
./venv/bin/python api_real.py
\end{verbatim}

\begin{warningbox}[\faExclamationTriangle\ Importante]
Siempre detenga el servicio API antes de restaurar la base de datos para evitar corrupción de datos.
\end{warningbox}

\subsection{Endpoints principales de la API}
\begin{longtable}{l l p{6cm}}
\toprule
\textbf{Método} & \textbf{Endpoint} & \textbf{Descripción} \\
\midrule
\endhead

POST & \texttt{/api/auth/login} & Autenticación de usuario \\
GET & \texttt{/api/auth/me} & Información del usuario autenticado \\
GET & \texttt{/api/usuarios} & Listar todos los usuarios \\
POST & \texttt{/api/usuarios} & Crear nuevo usuario \\
PUT & \texttt{/api/usuarios/:id} & Actualizar usuario \\
DELETE & \texttt{/api/usuarios/:id} & Eliminar usuario \\
GET & \texttt{/api/sentiment/distribution} & Distribución de sentimientos \\
GET & \texttt{/api/sentiment/kpis} & KPIs de sentimiento \\
GET & \texttt{/api/sentiment/wordcloud} & Nube de palabras \\
GET & \texttt{/api/ai/alerts} & Listar alertas \\
GET & \texttt{/api/ai/alerts/stats} & Estadísticas de alertas \\
PUT & \texttt{/api/ai/alerts/:id/resolve} & Resolver alerta \\
GET & \texttt{/api/nlp/keywords} & Palabras clave TF-IDF \\
GET & \texttt{/api/nlp/topics} & Tópicos LDA \\
GET & \texttt{/api/nlp/clusters} & Clusters K-Means \\
GET & \texttt{/api/osint/summary} & Resumen OSINT \\
POST & \texttt{/api/reports/generate} & Generar reporte \\
GET & \texttt{/api/configuracion-alertas} & Configuración de alertas \\
GET & \texttt{/api/logs} & Logs de actividad \\

\bottomrule
\caption{Endpoints principales de la API REST}
\end{longtable}

% ──────────────────────────────────────────────────────────────────────────
\section{Glosario}

\begin{description}[style=nextline, leftmargin=3cm]
  \item[OSINT] \textit{Open Source Intelligence} — Inteligencia de fuentes abiertas.
  \item[NLP] \textit{Natural Language Processing} — Procesamiento de Lenguaje Natural.
  \item[ML] \textit{Machine Learning} — Aprendizaje automático.
  \item[TF-IDF] \textit{Term Frequency -- Inverse Document Frequency} — Algoritmo de ponderación de palabras.
  \item[LDA] \textit{Latent Dirichlet Allocation} — Modelo generativo para extracción de tópicos.
  \item[K-Means] Algoritmo de clustering por partición en $k$ grupos.
  \item[NER] \textit{Named Entity Recognition} — Reconocimiento de entidades nombradas.
  \item[KPI] \textit{Key Performance Indicator} — Indicador Clave de Rendimiento.
  \item[UEBU] Unidad de Evaluación y Bienestar Universitario.
  \item[SADUTO] Sistema de Analítica de Datos Utilizando Técnicas Open Source Intelligence — Caso: Vicerrectorado de Grado.
  \item[ISG] Índice de Satisfacción General.
  \item[CRUD] \textit{Create, Read, Update, Delete} — Operaciones básicas sobre datos.
  \item[API REST] Interfaz de programación basada en arquitectura REST.
  \item[JWT] \textit{JSON Web Token} — Estándar de autenticación basado en tokens.
\end{description}

% ══════════════════════════════════════════════════════════════════════════
\newpage
\section*{Control de versiones del documento}
\begin{table}[H]
\centering
\begin{tabularx}{\textwidth}{c c l X}
\toprule
\textbf{Versión} & \textbf{Fecha} & \textbf{Autor} & \textbf{Cambios} \\
\midrule
1.0 & Feb 2026 & Alvaro Santiago Encinas Flores & Versión inicial del manual de usuario. \\
\bottomrule
\end{tabularx}
\end{table}

\end{document}
